\section{自主 Agent 的设计}

\subsection{Agent 的定义与工作机制}

当 LLM 在关键能力上日趋成熟——理解复杂输入、进行推理与规划、可靠地使用工具、从错误中恢复——Agent 开始在生产环境中出现。

Agent 的工作流程如下：
\begin{enumerate}
    \item 从人类用户的命令或交互讨论开始。
    \item 任务明确后，Agent \textbf{独立}进行规划和操作。
    \item 在执行过程中，通过工具调用结果、代码执行等从环境获取「\textbf{Ground Truth}」来评估进展。
    \item 可以在检查点（Checkpoint）暂停，请求人类反馈。
    \item 任务完成时终止，也通常设有最大迭代次数等停止条件。
\end{enumerate}

\begin{importantbox}{Agent 的实现本质}
Agent 听起来复杂，但其实现往往非常直接：\textbf{本质上就是 LLM 在一个循环中基于环境反馈使用工具}。关键不在于复杂的架构，而在于精心设计的工具集和清晰的工具文档。工具的设计质量决定了 Agent 的上限。
\end{importantbox}

\subsection{何时使用 Agent}

Agent 适用于以下场景：
\begin{itemize}
    \item \textbf{开放性问题}：无法预测所需步骤数。
    \item \textbf{无法硬编码固定路径}：需要模型自主决策。
    \item \textbf{可信环境}：对模型的决策有一定信任度。
    \item \textbf{规模化需求}：需要在受信任环境中大规模执行任务。
\end{itemize}

\begin{warningbox}{Agent 的风险与代价}
Agent 的自主性意味着：(1) \textbf{更高的成本}——多轮工具调用的 token 消耗远超单次调用；(2) \textbf{错误可能复合}——一个步骤的错误可能在后续步骤中放大。Anthropic 建议在\textbf{沙箱环境}中进行充分测试，并设置适当的 Guardrails（如最大迭代次数、敏感操作需人工确认等）。
\end{warningbox}

\subsection{Agent 的实践案例}

Anthropic 自身的两个 Agent 实现：

\begin{table}[H]
\centering
\begin{tabular}{>{\bfseries}p{3.5cm} p{9cm}}
\toprule
案例 & 说明 \\
\midrule
SWE-bench Coding Agent & 基于 Pull Request 描述，自动编辑多个文件以解决真实 GitHub Issue。能在 SWE-bench Verified 基准测试中解决实际问题。 \\
Computer Use Agent & Claude 使用计算机（鼠标、键盘、屏幕）来完成各种任务，是多模态 Agent 的典型案例。 \\
\bottomrule
\end{tabular}
\caption{Anthropic 的 Agent 实践案例}
\end{table}

\subsection{本章小结}
Agent 是 Agentic 系统复杂度谱系的最高端，适合开放性、不可预知路径的任务。其实现本质是 LLM + 工具调用循环，关键在于工具设计而非架构复杂度。使用时需注意成本、错误复合风险，以及充分的沙箱测试。


